\section{Half-plane arches and interior visits}\label{c7:boundary}

The boundary argument starts from the local balance and fusion identities of
Section~\ref{j:local}. An added cylinder row measures the deficit between
the boundary balance and the arches that embed in that cylinder. We express
its derivative by a positive integral, estimate that integral, and then use
a wedge comparison to control the radius of an arch. The last subsection
relates boundary chords visiting an interior vertex to the planar nesting
partition.

\begin{proposition}[Boundary arch mass]\label{j:boundary}
Fix a port on a horizontal tessellation half-plane boundary. Let $W_R$
be the total critical weight $x^N$ of first-exit same-boundary arches
reaching Euclidean distance at least $R$ from that port. Then
$W_R=R^{-1/4+o(1)}$. If $H(d)$ is the mass of arches ending $d$ boundary
steps to the right, then
\[
 2\sin(\pi/8)\sum_{d>0}H(d)=1,\qquad
 \sum_{d\ge n}H(d)\le n^{-1/4+o(1)}.
\]
The constants are unchanged by lattice translations and rotations.
\end{proposition}

\subsection{The fused cylinder increment}

On a finite cylinder of circumference $M$ with closed rims, let $Z$ be the
partition of exactly one unoriented noncontractible polygon. As in
Section~\ref{j:local}, it is the coefficient of $n$ in the loop sum with
contractible-loop weight zero and noncontractible-loop weight
$n=p+p^{-1}$. Write $\mathcal T_p(\phi)$ for the oriented row transfer
matrix with row parameter $\phi$, and put $\alpha=\pi/8$.

The row matrices with the same twist commute. Indeed, multiply the local
exchange identity across the sites and trace over the two auxiliary spaces.
Their twist $p^{s_1+s_2}$ commutes with the exchanging matrix. This proves
commutation where that matrix is invertible, and then identically by
continuation. In particular, a stack matrix element between empty rims, and
its coefficient $Z$, is symmetric in its rows.

We will also need the order of the error in fusing two rows. Let
$U=\operatorname{im}J\subset V_1\otimes V_2$ for the embedding at
$2\alpha$ in Section~\ref{j:local}, and at site $j$ set
\[
 X_j(\phi)=R_{2j}(\phi)R_{1j}(\phi+2\alpha).
\]
The local intertwining preserves $U\otimes V_j$ and, through the
isomorphism $J:V\to U$, induces $R(\phi+\alpha)$. At $\phi=0$,
\[
 X_j(0)=P_{2j}R_{1j}(2\alpha),\qquad
 \operatorname{im}X_j(0)\subset U\otimes V_j,
\]
where $P_{2j}=R_{2j}(0)$ is the swap. The inclusion follows because
$R_{1j}(2\alpha)=J_{1j}J_{1j}^*$ and the swap carries the embedded spaces
$1,j$ to $1,2$. Thus the operator induced by $X_j(\phi)$ on the quotient
by $U\otimes V_j$ is $O(\phi)$. Its product over the $M$ sites is
$O(\phi^M)$.

The twist preserves this invariant subspace and satisfies
$p^{s_1+s_2}J=Jp^s$. Taking the auxiliary trace on the invariant subspace
is therefore the trace of the fused row; the similarity through $J$
introduces no factor involving $J^*J$. Taking the trace on the quotient
gives
\begin{equation}\label{j:boundary-row-fusion}
 \mathcal T_p(\phi)\mathcal T_p(\phi+2\alpha)
 =\mathcal T_p(\phi+\alpha)+O(\phi^M).
\end{equation}
This is a locally holomorphic statement in $(\phi,p)$ near $(0,i)$.
After taking an empty-rim stack matrix element, the loop sum is a
polynomial in $n=p+p^{-1}$. Since $dn/dp=2$ at $p=i$, extracting its
coefficient of $n$ preserves the order of vanishing.

Put one new outer row on $L$ honeycomb rows, whose parameters are all
$\alpha$, and subtract the partition without the new row. We next describe
the height limit of this increment. A loop wholly in the new row gives
$v^M$. Otherwise the lift of the loop uses $2k>0$ ports into the base strip
per period. No exterior port is used. The new row therefore pairs
consecutive ports in $k$ occupied intervals, each with one turn of weight
$u$, one of weight $d$, and straight tiles of weight $v$.

Here the topology forces a particularly simple matching in the base.
Number the lifted ports in increasing order so that the outer intervals
are $(2j-1,2j)$. In a noncrossing lower perfect matching, the interior of
a pair is matched internally, so the pair's endpoints have opposite
parity. If its left endpoint is odd, the block including the endpoints is
closed under both matchings. If its left endpoint is even but the right
endpoint is not the next port, the interior is closed under both
matchings. Either case creates a finite component. The lift of one
noncontractible polygon is an infinite component, so the lower pairs are
exactly $(2j,2j+1)$.

For $1\le d<M$, let $H_M(d)$ be the weight of half-plane arches from a
fixed port to the port $d$ steps to its right that are disjoint from all
their nontrivial horizontal translates by $M$. More generally the
coefficient for $k$ lower arches requires the arches and all their
translates to be mutually disjoint. At fixed $M$ these coefficient sums
converge as $L\to\infty$: dropping the mutual disjointness conditions
bounds each by a product of the finitely many $H(d)$ involved. The coarse
summability proved in Section~\ref{j:local} suffices for this limit.

The tile weights satisfy
\[
 ud=c_\alpha v(1-v),\qquad c_\alpha=2\sin\alpha.
\]
There are at most $M-2k$ straight tiles in a $k$-interval term. Hence the
limiting increment is a polynomial $F(v)$ of degree at most $M$. For one
lower gap of length $d$, the upper interval has $M-d-1$ straight tiles
and $M$ possible positions in a labeled period. These positions count
each unoriented polygon once. Isolating these terms gives
\begin{equation}\label{j:boundary-arch-derivative}
 \begin{split}
 F(v)&=v^M+
 Mc_\alpha v(1-v)\sum_{d=1}^{M-1}v^{M-d-1}H_M(d)\\
 &\qquad +(1-v)^2Q_M(v),\\
 F(1)&=1,\qquad
 F'_v(1)=M\left(1-c_\alpha\sum_{d=1}^{M-1}H_M(d)\right),
 \end{split}
\end{equation}
for a polynomial $Q_M$ (zero when no term with two intervals is possible).

We now pass \eqref{j:boundary-row-fusion} to these increments. Commute its
two added rows to opposite closed rims of the base stack.
Inclusion--exclusion says that the difference between the two-row
increment and the sum of the one-row increments consists precisely of
loops touching both added rows. On a closed rim no tile can use a double
turn. On a fixed complex parameter neighborhood each remaining outer
weight is bounded in absolute value by a constant times the corresponding
positive honeycomb weight. There are only $2M$ outer tiles. In the arch
decomposition from the top, at least one base arch of a loop touching both
rims reaches depth comparable to $L$. Its weight tends to zero by
monotone convergence in the finite sums $H(d)$, after dropping the other
disjointness conditions. The same domination makes this convergence
locally uniform in the two outer parameters.

The one-row increments converge coefficientwise as polynomials in $v$,
and therefore locally uniformly in their parameters. Cauchy's formula now
passes every vanishing derivative of order below $M$ in
\eqref{j:boundary-row-fusion} to the limiting increments. With
\[
 y=i(2\phi-3\alpha),\qquad
 v=\frac{\cosh y-\sin\alpha}{\cosh y+\cos\alpha},\qquad
 f(y)=1-F(v),
\]
the resulting identity near $y=-3i\alpha$ is
\begin{equation}\label{j:boundary-fusion-jet}
 f(y)+f(y+4i\alpha)-f(y+2i\alpha)
 =1+O((y+3i\alpha)^M).
\end{equation}

\subsection{A positive moment representation}

For $k>0$, introduce column vectors indexed by $j=0,\ldots,M-1$ and
positive scalar weights
\[
 \begin{gathered}
 A_j(k)=\partial_\alpha^j\sinh(\alpha k),\qquad
 B_j(k)=\partial_\alpha^j\cosh(\alpha k),\\
 w(k)=\frac{2(1+4\sinh^2(\alpha k))}{\sinh(\pi k)},\qquad
 w_0(k)=k w(k).
 \end{gathered}
\]
Since $F(1)=1$, the function $f(y)$ is a linear combination of
$(\cosh y+\cos\alpha)^{-r}$, $1\le r\le M$. The first $M$ parameter
derivatives of
\[
 \frac{\sin\alpha}{\cosh y+\cos\alpha}
 =\int_{\mathbb R}e^{iky}\frac{\sinh(\alpha k)}{\sinh(\pi k)}\,dk
\]
span these powers: the derivative of order $j$ has a nonzero coefficient
of the denominator power $j+1$. Lemma~\ref{j:fourier-kernel} therefore
gives a vector $\mathbf c$ such that
\[
 f(y)=\int_{\mathbb R}e^{iky}
       \frac{A(k)^T\mathbf c}{\sinh(\pi k)}\,dk.
\]
The integral and its required derivatives converge locally uniformly in
$|\Im y|<\pi-\alpha$.

At $y=-3i\alpha$, the three arguments in
\eqref{j:boundary-fusion-jet} are $-3i\alpha,i\alpha,-i\alpha$, all in
this strip. Split the Fourier integral into positive and negative $k$.
For derivative order $j$, this gives a factor $2i^jk^j$ times
$\cosh(3\alpha k)$ if $j$ is even, and times
$\sinh(3\alpha k)-2\sinh(\alpha k)$ if $j$ is odd. Since $B_j(k)$ is
$k^j\cosh(\alpha k)$ or $k^j\sinh(\alpha k)$ in the two cases, the identities
\[
 \frac{\cosh(3z)}{\cosh z}
 =\frac{\sinh(3z)-2\sinh z}{\sinh z}
 =1+4\sinh^2z
\]
then turn the $M$ derivative equations into
\begin{equation}\label{j:boundary-moment}
 D\mathbf c=B(0),\qquad
 D=\int_0^\infty B(k)A(k)^T w(k)\,dk.
\end{equation}

Write $\Delta(z)=\prod_{i<j}(z_j-z_i)$ for the Vandermonde. The matrix
$D$ is nonsingular. Indeed, the ordered determinant expansion
gives
\[
 \begin{aligned}
 \det D={}&\int_{0<k_1<\cdots<k_M}
 \det[B(k_1),\ldots,B(k_M)]\\
 &\quad\cdot\det[A(k_1),\ldots,A(k_M)]
 \prod_iw(k_i)\,dk_i.
 \end{aligned}
\]
Both determinants in the integrand are strictly positive. To see this,
expand $\sinh(uk)$ or $\cosh(uk)$ in its positive power series. For each
increasing list of powers $p_\ell$, the coefficient minor after taking
successive derivatives at $u=\alpha$ is a positive factor times
\(\det[(p_\ell)_{\underline j}]=\Delta(p)\). The other minor is a
generalized Vandermonde in the ordered positive $k_i$, also positive.
Its nonvanishing follows by the usual zero bound obtained from Rolle's
theorem after successively dividing out the lowest power; its sign follows
by separating the positive $k_i$ in scale. The positive series expansion
now proves the claim.

Set
\[
 K(z)=A(z)^TD^{-1}B(0).
\]
For real $y\to+\infty$, the residue at $k=i$ in the Fourier formula for
$f$ gives $-2iK(i)e^{-y}$ as its first term. On the polynomial side,
\[
 1-v=2(\cos\alpha+\sin\alpha)e^{-y}+O(e^{-2y}).
\]
Equating the coefficients proves
\begin{equation}\label{j:boundary-residue}
 F'_v(1)=\frac{-iK(i)}{\cos\alpha+\sin\alpha}.
\end{equation}
We next express the quantity on the right by a positive expectation.

Give $0<k_1<\cdots<k_M$ the probability density proportional to
\begin{equation}\label{j:boundary-determinant-law}
 \begin{gathered}
 \det[A(k_1),\ldots,A(k_M)]\\
 {}\cdot\det[B(k_1),\ldots,B(k_M)]\prod_iw(k_i).
 \end{gathered}
\end{equation}
This law has a useful orbital form. For
\[
 D_k=\operatorname{diag}(k_1,\ldots,k_M,-k_1,\ldots,-k_M),
 \qquad
 Q=\mathcal U\operatorname{diag}(\alpha I_M,-\alpha I_M)\mathcal U^*,
\]
where $\mathcal U$ has normalized Haar measure on $U(2M)$, put
\[
 I(k)=\mathbb E\exp(\operatorname{tr}(D_kQ)).
\]
Then \eqref{j:boundary-determinant-law} is proportional to
\begin{equation}\label{j:boundary-orbital-law}
 \Delta(k^2)^2\prod_iw_0(k_i)\,I(k).
\end{equation}
In particular it is a positive probability density.

For completeness, we derive the paired confluent form used here from the
Harish-Chandra--Itzykson--Zuber integral
\cite[Section~2, (2.1)]{ZinnJustinZuber2003}. For distinct real lists
$z,u$ of size $d$, put $D_z=\operatorname{diag}(z)$ and
$D_u=\operatorname{diag}(u)$. Normalized unitary integration gives
\[
 \int_{U(d)}e^{\operatorname{tr}(D_z\mathcal U D_u\mathcal U^*)}\,d\mathcal U
 =C_d\frac{\det[e^{z_i u_j}]}{\Delta(z)\Delta(u)},
 \qquad C_d>0.
\]
One derivation uses the isotropic compression mechanism of
\cite[Section~4]{Baryshnikov2001}. Compress $\mathcal U D_u\mathcal U^*$ to a leading
$(d-1)$-dimensional minor. Its ordered eigenvalues $l_i$ interlace the
ordered $u_i$, and the last resolvent entry is
\[
 \sum_{j=1}^d\frac{\omega_j}{t-u_j}
 =\frac{\prod_{j<d}(t-l_j)}{\prod_{j\le d}(t-u_j)}.
\]
Here $\omega$ is uniform on the simplex, since its entries are squared
moduli of a uniform complex unit vector. The residues give a bijection
between the simplex interior and the interlacing region. Its absolute
Jacobian is $\Delta(l)/\Delta(u)$, as follows by differentiating the
residues, or by passing through the coefficients of the monic numerator.
The compressed matrix, conditional on its eigenvalues, is again
unitarily invariant. Induction with test parameters $z_i-z_d$ integrates
each exponential column over $[u_j,u_{j+1}]$; column subtraction and the
factor $e^{z_d\sum u_j}$ give the displayed determinant formula. The
simplex normalization changes only $C_d$.

Apply this formula to the lists $\pm k_i,\pm u_i$. The absolute paired
Vandermonde is
\[
 |\Delta(k_1,\ldots,k_M,-k_1,\ldots,-k_M)|
 =2^M\prod_i k_i\,\Delta(k^2)^2,
\]
and similarly for $u$. Splitting the exponential determinant into its
sinh and cosh blocks therefore gives a positive constant times
\[
 \frac{\det[\sinh(k_i u_j)]\det[\cosh(k_i u_j)]}
      {\Delta(k^2)^2\Delta(u^2)^2\prod_i k_i u_i}.
\]
Divide both numerator determinants by $\Delta(u)$ and let $u_i\to\alpha$
by divided differences. The remaining $u$ factors are constants, and
the numerator becomes the two derivative determinants in
\eqref{j:boundary-determinant-law}. This proves
\eqref{j:boundary-orbital-law}. Also
\[
 I(k)\le e^{2\alpha\sum_i k_i},\qquad
 w_0(k)\sim4k e^{-6\alpha k}\quad(k\to\infty),\qquad
 w_0(0)=2/\pi,
\]
so the density is integrable and extends continuously when one point
approaches zero.

\begin{lemma}[Pinning at the origin]\label{j:orbital-pin}
Let $\rho(0)$ be the one-point intensity at zero of
\eqref{j:boundary-orbital-law}. Pinning one point there leaves $N=M-1$
positive points with density proportional to
\[
 \Delta(k^2)^2\prod_{j=1}^N k_j^4w_0(k_j)\,I(k_1,\ldots,k_N,0).
\]
The orbital integral here still uses the $2M$-dimensional orbit, with one
zero pair in $D_k$. For its expectation $\mathbb E_{\rm pin}$,
\[
 \begin{aligned}
 \frac{\sin\alpha}{\alpha}\frac{\rho(0)}{w_0(0)}
   \mathbb E_{\rm pin}\prod_{j=1}^N(1+k_j^{-2})
 &\le -iK(i)\\
 &\le\frac{\rho(0)}{w_0(0)}
   \mathbb E_{\rm pin}\prod_{j=1}^N(1+k_j^{-2}).
 \end{aligned}
\]
Empty products have value one.
\end{lemma}

\begin{proof}
The continuous density at a single zero point defines $\rho(0)$.
For fixed $M$, the probability of two points in $[0,t]$ is $O_M(t^6)$
by their squared Vandermonde factor. Hence
\[
 \lim_{t\downarrow0}\frac{\mathbb P(\min_i k_i<t)}t=\rho(0).
\]
In the determinant density, evaluating a point at zero inserts
$A'(0)$ and $B(0)$, since $A(k)=kA'(0)+O(k^3)$ and
$w(k)=w_0(0)/k+O(k)$. Define
\[
 \mathcal R(k)=
 \frac{\det[A(i)/i,A(k_1),\ldots,A(k_N)]}
      {\det[A'(0),A(k_1),\ldots,A(k_N)]}.
\]
The ordered adjugate expansion for $K(i)=A(i)^TD^{-1}B(0)$ then gives
exactly
\[
 \frac{K(i)}i=\frac{\rho(0)}{w_0(0)}
 \mathbb E_{\rm pin}\mathcal R(k).
\]

To compare the two determinants, first take distinct
$0<u_1<\cdots<u_M$ in a fixed neighborhood of $\alpha$ contained in
$(0,\pi)$, before taking their confluent limit. Put
$f_z(u)=\sinh(zu)/z$, with $f_0(u)=u$; thus $f_i(u)=\sin u$.
In the determinant with first column $f_z(u)$ and remaining columns
$\sinh(k_j u)$, remove $\prod_i f_z(u_i)$ and subtract adjacent rows.
The identity
\[
 \left(\frac{\sinh(k_jt)}{f_z(t)}\right)'
 =\frac{k_j^2-z^2}{f_z(t)^2}
   \int_0^t f_z(v)\sinh(k_jv)\,dv
\]
follows by differentiating the numerator of the left side and using the
two second-order differential equations. After removing
$\prod_j(k_j^2-z^2)$, row differences on the inner integrals express the
remaining determinant as an integral of
\[
 \det[\sinh(k_jv_i)]\prod_i\frac{f_z(v_i)}{f_z(t_i)^2}
\]
over $u_i<t_i<u_{i+1}$ and $t_{i-1}<v_i<t_i$, with $t_0=0$.
Its measure is positive by the same generalized Vandermonde argument.

Let $g(t)=\sin t/t$ with $g(0)=1$. The ratio for $z=i$ and $z=0$,
after removing $\prod_j(1+k_j^{-2})$, is a positive weighted average of
\[
 R=\frac{\prod_{i=1}^M g(u_i)\prod_{i=1}^{M-1}g(v_i)}
         {\prod_{i=1}^{M-1}g(t_i)^2}.
\]
The function $\log g$ decreases on $(0,\pi)$. The two interlacings give
\[
 \log g(u_M)\le\sum_i\log g(u_i)-\sum_i\log g(t_i)
 \le\log g(u_1),
\]
and
\[
 0\le\sum_i\log g(v_i)-\sum_i\log g(t_i)
 \le-\log g(t_{M-1}).
\]
Consequently $g(u_M)\le R\le g(u_1)/g(u_M)$, uniformly in $M$ and the
$k_j$. Taking all $u_i\to\alpha$ bounds this residual ratio between
$\sin\alpha/\alpha$ and one. The exact adjugate formula and $1/i=-i$
prove the lemma.
\end{proof}

\subsection{Orbital derivatives and stochastic comparison}

We will apply the same estimates to the unpinned and pinned densities.
For $e=0$, write $\mathbb O_N^{(0)}$ for the unpinned orbital law of size
$N$. For $e=4$, write $\mathbb O_N^{(4)}$ for the law obtained by pinning
one point of the size-$N+1$ orbital law. Their densities have the form
\[
 \Delta(k^2)^2\prod_{i=1}^N k_i^e w_0(k_i)\,I^{(e)}(k),
\]
where $I^{(4)}$ retains its zero pair. All coordinates below are sorted
and positive.

\begin{lemma}[Orbital derivative bound]\label{j:orbital-ward}
For either orbital factor, there are numbers $q_j$ such that
\[
 \partial_{k_j}\log I^{(e)}=2q_j,\qquad
 0\le q_j\le
 \min\left\{\alpha,\alpha^2k_j,
 \frac{\alpha^2}{\sum_{\ell=1}^N(k_j+k_\ell)^{-1}}\right\}.
\]
The bounds are uniform in $N$.
\end{lemma}

\begin{proof}
Use brackets for expectation tilted by
$\exp(\operatorname{tr}(D_kQ))$, and set $q_j=\langle Q_{+j,+j}\rangle$.
Exchanging all paired positions and sending $Q$ to $-Q$ preserves the
orbit and this tilt. It gives
$\langle Q_{-j,-j}\rangle=-q_j$, proving the derivative formula.
The orbital integral is even in each $k_j$, while its logarithm is
convex in that coordinate as a logarithmic moment generating function.
Thus $q_j\ge0$. Also $q_j\le\alpha$ because $\|Q\|=\alpha$.

Infinitesimal conjugation mixing two indices $a,b$ gives
\[
 \langle Q_{aa}-Q_{bb}\rangle
 =((D_k)_{aa}-(D_k)_{bb})\langle |Q_{ab}|^2\rangle.
\]
For example, differentiate the integral of $\Re Q_{ab}$ times the tilt
under the generator $E_{ab}-E_{ba}$. Diagonal-unitary symmetry gives
$2\langle(\Re Q_{ab})^2\rangle=\langle|Q_{ab}|^2\rangle$, yielding the
displayed factor. For $a=+j,b=-\ell$, this becomes
\[
 \langle |Q_{+j,-\ell}|^2\rangle
 =\frac{q_j+q_\ell}{k_j+k_\ell}.
\]
The $\ell=j$ term and
$\sum_b\langle|Q_{+j,b}|^2\rangle=\alpha^2$, which follows from
$Q^2=\alpha^2\mathrm{Id}$, give $q_j\le\alpha^2k_j$. Sum over all
moving negative positions and drop $q_\ell\ge0$ to get
\[
 q_j\sum_{\ell=1}^N(k_j+k_\ell)^{-1}\le\alpha^2.
\]
In the pinned case the omitted zero positions only add nonnegative terms
to the full row sum.
\end{proof}

For a positive one-point weight $w_*$, let $\mathbb P_N[w_*]$ denote
the pure law with density $\Delta(k^2)^2\prod_iw_*(k_i)$, whenever
integrable. Write $\mu\preceq_{\rm st}\nu$ when the expectation of every
bounded coordinatewise increasing function under $\mu$ is at most its
expectation under $\nu$.

Every pure law used here is positively associated: the covariance of two
bounded increasing functions is nonnegative. Here is an induction that
also covers any fixed upper endpoint. The size-one assertion is the
one-dimensional covariance inequality. Given the last point $t$, the
earlier points have the same pure form with weight
$w_*(k)(t^2-k^2)^2$ below $t$. If $s<t$, the likelihood ratio for
conditioning at $s$ relative to conditioning at $t$ is proportional to
\[
 {\bf1}_{k_{N-1}<s}\prod_{i<N}
 \left(\frac{s^2-k_i^2}{t^2-k_i^2}\right)^2.
\]
It decreases in every coordinate, including across its cutoff.
Association at size $N-1$ therefore makes the earlier points
stochastically increasing in the last point. Conditional covariance and
the one-dimensional covariance inequality then prove association at size
$N$. In particular an increasing (respectively decreasing) likelihood
ratio tilting a pure law moves the points upward (respectively downward)
in stochastic order. Truncation extends this rule to the unbounded
nonnegative tests used below.

The logarithmic derivative of $w_0$ is
\[
 D_0(k)=\frac1k+3\alpha\tanh(3\alpha k)
        -\alpha\tanh(\alpha k)-8\alpha\coth(8\alpha k).
\]
It satisfies
\begin{equation}\label{j:boundary-weight-slope}
 \begin{gathered}
 -6\alpha<D_0(k)<0,\qquad D_0(k)\longrightarrow-6\alpha,\\
 D_0(k)=-\tfrac{40}{3}\alpha^2 k+O(k^3)\quad(k\downarrow0).
 \end{gathered}
\end{equation}
For the upper bound put $u=\alpha k$ and use
$8\coth(8u)>1/u+4\tanh(4u)$. For the lower bound, it suffices to use
\[
 \frac{16u}{e^{16u}-1}+\frac{6u}{e^{6u}+1}<1,
\]
which follows from $e^{16u}\ge1+16u+(16u)^2/2$ and
$e^{6u}\ge1+6u$. The limits follow by expansion.

By Lemma~\ref{j:orbital-ward}, the logarithmic derivative of the
orbital density relative to the pure density with weight
$k^e w_0(k)e^{2\alpha k}$ is $2q_i-2\alpha\le0$. Hence
\begin{equation}\label{j:boundary-upper-pure}
 \mathbb O_N^{(e)}
 \preceq_{\rm st}\mathbb P_N[k^e w_0(k)e^{2\alpha k}],
 \qquad e=0,4.
\end{equation}
For this pure law its first $j$ points are in turn stochastically below
$\mathbb P_j[k^e w_0(k)e^{2\alpha k}]$. Indeed, conditioning on the
later points multiplies the latter density by
$\prod_{i\le j<r}(k_r^2-k_i^2)^2$ and an upper cutoff, both decreasing
in the first $j$ coordinates.

The largest point of this standalone size-$j$ law is at most $Cj$ with
probability $1-O(e^{-cj})$, with constants valid for $e=0,4$. To verify
the needed uniformity, condition on all but one of its unsorted points.
The remaining density is proportional to $p(k)^2W(k)$, where
$\deg p\le2j$ and $W(k)=k^e w_0(k)e^{2\alpha k}$. On $[j,2j]$,
$W\ge c e^{-C_0j}$, while $W$ is bounded above by a fixed polynomial
times an exponentially decaying function. Choose $2j+1$ interpolation
points in uniformly spaced small subintervals of $[j,2j]$, with widths
and intervening gaps bounded below, at which
$|p|\le C\|p\|_{L^2[j,2j]}$. Lagrange interpolation, using denominators
bounded below by $c^{2j}l!(2j-l)!$, gives for $k>2j$
\[
 |p(k)|\le\|p\|_{L^2[j,2j]}
       [C_1(k+2j)/j]^{2j}.
\]
After weakening the exponential decay to absorb the fixed polynomial,
this implies, for a sufficiently large fixed $A$,
\[
 \begin{aligned}
 \frac{\int_{Aj}^\infty p(k)^2W(k)\,dk}
      {\int_0^\infty p(k)^2W(k)\,dk}
 &\le C e^{C_0j}\int_{Aj}^\infty
       [C_1(k+2j)/j]^{4j}e^{-c_0k}\,dk
 \\ &\le C e^{-c_1j}.
 \end{aligned}
\]
The last inequality follows after $k=jt$ because the linear negative
term in $t$ eventually dominates $4\log(C_1(t+2))$. A union bound proves
the maximum estimate.

Fix $0<\epsilon<1$ and let $m=\lfloor N^{1-\epsilon}\rfloor$. Applying
the preceding bounds to the first $j$ points and summing over $j\ge m$
shows that under either orbital law the event
\[
 E=\{k_j\le C\max(m,j)\text{ for every }j\}
\]
has probability $1-O(e^{-c m})$. On $E$, for $i\le m$,
\[
 \sum_{\ell=1}^N(k_i+k_\ell)^{-1}
 \ge c\sum_{\ell=m}^N(m+\ell)^{-1}
 \ge c'\log(N/m).
\]
Lemma~\ref{j:orbital-ward} therefore gives
\begin{equation}\label{j:boundary-small-orbit}
 \sup_{\substack{i\le m\\ k\in E}}q_i
 \le\frac{C}{\epsilon\log N}=o(1)
\end{equation}
for each fixed $\epsilon$.

\subsection{Gamma comparison and evaluation}

For $\ell,\beta>0$ define
\[
 W_{\ell,\beta}(k)=|\Gamma(\ell+i\beta k)|^2,\qquad b=\pi\beta.
\]
These weights will bound the points in the direction needed for the
decreasing product in Lemma~\ref{j:orbital-pin}. For the event $E$ above,
write $k_{\le m}=(k_1,\ldots,k_m)$ and
$k_{>m}=(k_{m+1},\ldots,k_N)$.

\begin{lemma}[Gamma comparison]\label{j:boundary-gamma-comparison}
For every fixed $b_+>6\alpha$, there is $\ell_+>0$ such that, with
$\beta_+=b_+/\pi$, for all $N$ and $e=0,4$,
\[
 \mathbb P_N[k^e W_{\ell_+,\beta_+}]
 \preceq_{\rm st}\mathbb O_N^{(e)}.
\]
For fixed $0<b_-<6\alpha$ and $0<\epsilon<1$, there is $\ell_->0$ such
that, with $\beta_-=b_-/\pi$ and for sufficiently large $N$,
\[
 \mathcal L_{\mathbb O_N^{(e)}}(k_{\le m}\mid E,k_{>m})
 \preceq_{\rm st}\mathbb P_m[k^eW_{\ell_-,\beta_-}]
\]
for almost every later configuration compatible with $E$, where
$m=\lfloor N^{1-\epsilon}\rfloor$.
\end{lemma}

\begin{proof}
The Euler product gives
\[
 G_{\ell,\beta}(k):=(\log W_{\ell,\beta})'(k)
 =-2\beta^2k\sum_{n\ge0}
       \frac1{(n+\ell)^2+\beta^2k^2}.
\]
The product for the ratio to $\Gamma(\ell)^2$ can also be obtained by
approximating the Gamma integral by the beta integral, using
$(1-y/L)^L\to e^{-y}$ and integration by parts. Integral comparison of
the decreasing summand gives
\[
 \begin{array}{ll}
 G_{\ell,\beta}(k)\le-b+2/k,&0<\ell\le1,\\[2mm]
 G_{\ell,\beta}(k)\ge-b,\quad
 |G_{\ell,\beta}(k)|/k\le2\beta^2(\ell^{-2}+\ell^{-1}),&\ell\ge1.
 \end{array}
\]
The first summand also gives
\[
 G_{\ell,\beta}(k)/k
 \le-\frac{2\beta^2}{\ell^2+\beta^2k^2}.
\]

Suppose first that $b>6\alpha$. At fixed $k>0$, as $\ell\downarrow0$,
the sum dominates its integral strictly, so
$G_{\ell,\beta}(k)$ tends to a value below $-b<-6\alpha<D_0(k)$.
On a fixed compact interval away from zero this inequality holds uniformly
for sufficiently small $\ell$. In a sufficiently small interval at zero,
the last displayed bound and the boundedness of $D_0(k)/k$ give the same
comparison. In the tail use $G_{\ell,\beta}\le-b+2/k$ and
$D_0>-6\alpha$. These three regions select one fixed $\ell>0$ with
$G_{\ell,\beta}\le D_0$ everywhere. The likelihood ratio of the orbital
law to $\mathbb P_N[k^eW_{\ell,\beta}]$ is then increasing, since its
logarithmic derivative is $D_0+2q_i-G_{\ell,\beta}\ge0$. Association of
the pure law proves the first assertion.

Now fix $b<6\alpha$ and $\epsilon>0$. Near zero,
\eqref{j:boundary-weight-slope} and $q_i\le\alpha^2k_i$ give
\[
 D_0(k_i)+2q_i
 \le-\bigl(\tfrac{34}{3}+o(1)\bigr)\alpha^2k_i.
\]
Choose $c_0>0$ and a fixed small interval on which
$D_0(k)+2\alpha^2k\le-c_0k$. Then choose $\ell$ so large that
$2\beta^2(\ell^{-2}+\ell^{-1})<c_0$. The bound on the Gamma derivative
gives $G_{\ell,\beta}(k)\ge-c_0k\ge D_0(k)+2q_i$ on that interval.
On a fixed compact
interval away from zero, $D_0<0$, the Gamma derivative tends uniformly
to zero as $\ell$ grows, and \eqref{j:boundary-small-orbit} makes
$q_i=o(1)$ uniformly on $E$. In the tail, use
$D_0\to-6\alpha<-b$, $G_{\ell,\beta}\ge-b$, and the same bound on $q_i$.
Thus, for one fixed large $\ell$ and sufficiently large $N$,
\[
 G_{\ell,\beta}(k_i)\ge D_0(k_i)+2q_i
 \quad\text{for }i\le m\text{ on }E.
\]

Given the later points and $E$, the likelihood ratio of the first $m$
points to $\mathbb P_m[k^eW_{\ell,\beta}]$ is proportional to
\[
 {\bf1}_{k_m<\min(Cm,k_{m+1})}
 I^{(e)}(k)\prod_{i\le m}\frac{w_0(k_i)}{W_{\ell,\beta}(k_i)}
 \prod_{\substack{i\le m\\r>m}}(k_r^2-k_i^2)^2.
\]
The factors $k_i^e$ have canceled. Inside this downward closed support,
its logarithmic derivative in $k_i$ is
\[
 D_0(k_i)+2q_i-G_{\ell,\beta}(k_i)
 -4k_i\sum_{r>m}(k_r^2-k_i^2)^{-1}\le0.
\]
It is therefore a decreasing ratio, including across the cutoff.
Association of the pure comparison law proves the conditional assertion.
\end{proof}

We next compute the particular pure statistic appearing here. Throughout
this calculation $\ell,\beta>0$ are fixed. Normalize
$|\Gamma(\ell+iu)|^2\,du$ to a probability measure on $\mathbb R$.
Its Fourier transform is $\cosh(t/2)^{-2\ell}$: writing
\[
 \Gamma(\ell+iu)=\int_{\mathbb R}e^{(\ell+iu)y-e^y}\,dy
\]
reduces the transform of its squared modulus to the autocorrelation of
$e^{\ell y-e^y}$, which is evaluated by scaling $e^y$. Shifting this
Fourier contour by a small positive or negative imaginary amount also
gives exponential decay of the weight. Thus the generating function
\[
 (1-iz)^{-\ell+iu}(1+iz)^{-\ell-iu}
 =\sum_{n\ge0}P_n(u)z^n
\]
gives the Meixner--Pollaczek polynomials
$P_n^{(\ell)}(u;\pi/2)$ \cite[Equation~18.23.7]{DLMF}. Integrating the
product of two generating functions near zero gives
$(1-zz')^{-2\ell}$. Consequently the $P_n$ are real, have degree $n$
with leading coefficient $2^n/n!$, have parity $n$, and have squared norms
\[
 c_n=\frac{(2\ell)_n}{n!}\sim\frac{n^{2\ell-1}}{\Gamma(2\ell)}.
\]
Their evaluations needed below satisfy
\[
 \begin{aligned}
 (-1)^jP_{2j}(0)&=\frac{(\ell)_j}{j!}
 \sim\frac{j^{\ell-1}}{\Gamma(\ell)},\\
 (-1)^jP_{2j}(i\beta)&
 \sim\frac{2^{2\beta-1}j^{\ell+\beta-1}}{\Gamma(\ell+\beta)}.
 \end{aligned}
\]
For the second equivalent, take the even coefficients of
$(1-t)^{-\ell-\beta}(1+t)^{-\ell+\beta}$. If $\ell\ge\beta$, factor it
as $(1-t^2)^{-\ell+\beta}(1-t)^{-2\beta}$ and use the positive binomial
convolution. For $\ell>\beta$, its binomial equivalents, with the
summation index divided by $j$, give the leading constant
\[
 \frac{2^{2\beta-1}}{\Gamma(\ell-\beta)\Gamma(2\beta)}
 \int_0^1v^{\ell-\beta-1}(1-v)^{2\beta-1}\,dv
 =\frac{2^{2\beta-1}}{\Gamma(\ell+\beta)};
\]
the case $\ell=\beta$ is immediate. If $\ell<\beta$, the coefficients of
$(1+t)^{\beta-\ell}$ are $O((1+n)^{\ell-\beta-1})$ and are absolutely
summable. Convolution with the first factor is asymptotic to its
coefficient times their sum $2^{\beta-\ell}$; the indices of the
summable factor above $n/2$ contribute only $O(n^{2\ell-1})$, smaller
than the main term. This proves the displayed equivalents.

Use the even polynomials $P_{2j}(\beta k)$ on the positive axis.
Their evaluation kernel through $j=L-1$ is the kernel of
$\mathbb P_L[W_{\ell,\beta}]$, up to its fixed one-point normalization.
The determinant expansion for its intensity at zero therefore gives
\[
 \rho_{\Gamma,L}(0)\asymp_{\ell,\beta}
 \sum_{j<L}\frac{P_{2j}(0)^2}{c_{2j}}
 \asymp_{\ell,\beta}1+\log L.
\]
Pinning one point of this law gives
$\mathbb P_{L-1}[k^4W_{\ell,\beta}]$; denote its expectation by
$\mathbb E_{\Gamma,L}^{\rm pin}$. Replacing evaluation at zero in
one determinant by evaluation at $i$ changes its Vandermonde by exactly
$\prod_j(1+k_j^{-2})$. Hence
\begin{equation}\label{j:boundary-gamma-statistic}
 \begin{split}
 \mathbb E_{\Gamma,L}^{\rm pin}\prod_j(1+k_j^{-2})
 &=\frac{\sum_{j<L}P_{2j}(0)P_{2j}(i\beta)/c_{2j}}
         {\sum_{j<L}P_{2j}(0)^2/c_{2j}}\\
 &\asymp_{\ell,\beta}\frac{L^\beta}{1+\log L}
 =L^{\beta+o(1)}.
 \end{split}
\end{equation}
Indeed the denominator summands are asymptotic to a positive constant
times $j^{-1}$, and the numerator summands to a positive constant times
$j^{\beta-1}$. In particular the exponent statement is for fixed
$\ell,\beta$.

We can now finish the estimate of \eqref{j:boundary-residue}.
For the unpinned orbital law, its lower stochastic bound in
Lemma~\ref{j:boundary-gamma-comparison} bounds its decreasing
minimum-tail event from above by the Gamma minimum tail. Conversely,
\eqref{j:boundary-upper-pure} and the comparison of the first point to
the standalone one-point pure law bound that event from below by the
one-point law with weight $w_0e^{2\alpha k}$. Taking $t\downarrow0$ for
each $M$ gives
\[
 c\le\rho(0)\le C(1+\log M).
\]
The constants may depend on the fixed comparison parameters.

For the pinned law, the product in Lemma~\ref{j:orbital-pin} is
decreasing. The first Gamma comparison therefore bounds its expectation
above by \eqref{j:boundary-gamma-statistic} with $L=N+1=M$, since the
pinned law has $N=M-1$ moving points. For a lower
bound, restrict to $E$ and discard all factors after $m$, which are at
least one. Conditional on the later points, the second Gamma comparison
bounds the remaining decreasing product below by the pure size-$m$
expectation, namely \eqref{j:boundary-gamma-statistic} with $L=m+1$.
Its bound is multiplied by $\mathbb P(E)=1-O(e^{-cm})$.

To record the order of limits, first fix $b_+>6\alpha$ for the upper
bound, and fix $b_-<6\alpha$ and $\epsilon>0$ for the lower bound.
Choose their $\ell$ parameters as in the comparison lemma, and only then
let $M\to\infty$. Put $a_M=\log F'_v(1)/\log M$. The pin and residue
formulas give
\[
 (1-\epsilon)\frac{b_-}{\pi}
 \le\liminf_{M\to\infty}a_M
 \le\limsup_{M\to\infty}a_M
 \le\frac{b_+}{\pi}.
\]
Letting these fixed parameters approach $6\alpha$ and then
$\epsilon\downarrow0$ proves
\begin{equation}\label{j:boundary-increment-exponent}
 F'_v(1)=M^{3/4+o(1)}.
\end{equation}

\subsection{Wedge comparison}\label{j:wedge-comparison}

We return to the physical arches. Define the positive deficit
\[
 \delta_M=\frac1{c_\alpha}-\sum_{d=1}^{M-1}H_M(d)
          =\frac{F'_v(1)}{c_\alpha M}
          =M^{-1/4+o(1)}.
\]
The coarse bound $c_\alpha\sum_{d>0}H(d)\le1$ from
Section~\ref{j:local}, together with $H_M(d)\le H(d)$ and
$\delta_M\to0$, gives the saturation
\[
 \sum_{d>0}H(d)=\frac1{c_\alpha}.
\]
Consequently, with $A(n)=\sum_{d\ge n}H(d)$,
\[
 A(M)\le\delta_M.
\]
Every arch contained within distance strictly less than $M/2$ of its
source is disjoint from its nontrivial translates by $M$. Thus the
arches counted in the deficit all reach distance at least $M/2$, giving
\[
 \delta_M\le W_{M/2}.
\]
Choosing $M$ just above $2R$ proves the required lower exponent for
$W_R$ and the stated upper exponent for $A(n)$.

For the upper radius bound, intersect the upper half-plane with the wedge
between rays of angles $0$ and $\theta=\pi/3$. A source on its horizontal
ray has distance $a\in\{1/2,3/2,\ldots\}$ from the vertex. The signed
balance holds on the infinite wedge. To justify this, first exhaust the
half-plane by convex tessellation polygons with extending horizontal
bases. The real contribution of their same-side arches tends to one by
saturation. Every other real boundary coefficient is at least
$\sin\alpha$, so the total raw weight of remote exits tends to zero.
Remote exits in the wedge intersections are a subset of these exits.
The fixed-ray exit sums converge absolutely by the same positive bound,
and the finite signed identities therefore pass to the wedge.

Put $s=3/8$. For same-side walks the phases are $e^{is\pi}$ to the left
and $e^{-is\pi}$ to the right, while exits to the other ray have phase
$e^{is\theta}$. Let $L_<(a)$ and $L_>(a)$ be the lost left and right
same-side weights on restriction from the half-plane to the wedge,
including endpoints removed by the restriction. Subtract the balances
and project perpendicular to $e^{is\theta}$. This gives
\[
 \sin(s(\pi-\theta))L_<(a)
 -\sin(s(\pi+\theta))L_>(a)=0,
 \qquad L_<(a)=rL_>(a),\quad r=\sqrt2.
\]

For positive half-integer ports $a,b$, let $D(a,b)$ be the lost
same-side weight with those fixed endpoints. Reversal gives
$D(a,b)=D(b,a)$. The missing negative ports to the left contribute
$A(\lceil a\rceil)$, so
\[
 L_<(a)=A(\lceil a\rceil)+\sum_{0<b<a}D(a,b),
 \qquad
 L_>(a)=\sum_{b>a}D(a,b).
\]
Summing $L_<=rL_>$ over $0<a\le T$ gives the positive identity
\[
 \begin{aligned}
 (r-1)\sum_{0<a<b\le T}D(a,b)
 +r\sum_{0<a\le T<b}D(a,b)
 &=\sum_{0<a\le T}A(\lceil a\rceil)\\
 &\le T^{3/4+o(1)}.
 \end{aligned}
\]
Because $r>1$, this bounds
$\sum_{0<a\le T}(L_<(a)+L_>(a))$ by the same order. After translating
the source to a fixed port, the wedges increase with $a$; their total
loss therefore decreases. The summed bound implies
\[
 L_<(a)+L_>(a)\le a^{-1/4+o(1)}.
\]
Intersect two reflected wedges whose vertices are horizontal distance
$a$ on opposite sides of the source. Their intersection is a triangle
of diameter $O(a)$. An arch reaching distance $Ca$ for sufficiently
large fixed $C$ leaves at least one wedge, so its weight is bounded by
the sum of these two losses. This proves the upper exponent for $W_R$
and completes Proposition~\ref{j:boundary}.

\subsection{A bound for segments between height extrema}

We record a segment bound used later in the marked-bond summability
argument. Measure vertex height in units of one third of the spacing
between horizontal tessellation lines. Heights are integers and change
by one or two on an edge.

\begin{lemma}\label{j:extremal-segment}
Fix a honeycomb vertex and two distinct height levels, with the fixed
vertex on one level. The total critical weight of SAW segments from
that vertex to any vertex on the other level, whose endpoints are
opposite height extrema of the segment, is bounded by an absolute
constant independent of the levels.
\end{lemma}

\begin{proof}
Extend each endpoint to a port on the adjacent horizontal line beyond
its extremal level. This requires at most one additional triangle at
each end, strictly beyond the segment's range of heights. Between two
successive horizontal lines the lower triangles have a port on the
bottom and the higher triangles a port on the top; a triangle of either
type has a neighbor of the other type. These observations provide the
extension without intersecting the segment. Choose an extension by a
fixed local rule. Trimming the vertices beyond the extremal levels
recovers the segment, and the possible choices have bounded multiplicity.

The extended walk is a first-exit crossing of the slab between the two
chosen lines. The convex boundary balance, exhausted horizontally,
bounds its total weight for each of the finitely many possible starting
ports. Removing the at most two extra vertices costs at most $x^{-2}$.
This gives a constant independent of the height difference.
\end{proof}

For the later use, split an $n$-edge walk at a highest vertex and treat
the two halves from that vertex. Within either half, go to the last
occurrence of its lowest remaining level, then to the last occurrence
of its highest remaining level, and continue alternately. These are
segments of Lemma~\ref{j:extremal-segment}. Their positive integer spans
decrease strictly after a possible equality of the first two, since each
chosen extremum was the last occurrence of its level. The sum of the
spans is at most $2n$, so there are $O(\sqrt n)$ segments and at most
$(2n+1)^{O(\sqrt n)}$ span lists. For a prescribed list, successive use
of the lemma bounds the endpoint-summed product of segment weights by a
constant to the number of segments. Recovering the weight of the joined
walk costs $x^{-1}$ per shared endpoint. These are the subexponential
concatenation bounds needed in Section~\ref{c7:marked}.

They also recover the reciprocal exponential growth rate of full-plane
SAWs. Let $c_n$ count $n$-edge SAWs from a fixed vertex. A highest vertex
has only $O(n^2)$ possible positions, and the split has at most $n+1$
positions. Summing the span lists and the preceding product bounds gives
$c_nx^{n+1}\le\exp(O(\sqrt n\log n))$, hence
$\limsup c_n^{1/n}\le x^{-1}$. Lattice symmetry gives
$c_{n+m}\le c_nc_m$, so the limit $\mu=\lim c_n^{1/n}$ exists. If
$\mu<x^{-1}$, then $c_nx^n\le Cr^n$ for some $r<1$. An arch reaching
radius $R$ has at least a constant times $R$ vertices, so summing this
bound over its possible lengths would make $W_R$ exponentially small,
contrary to Proposition~\ref{j:boundary}. Thus $\mu=x^{-1}$.
\cite{DuminilCopinSmirnov2012}

\subsection{Boundary chords and interior visits}

We first convert the arch estimate into the partition of macroscopic
chords. The following consequence of the signed balance makes the
comparison precise.

\begin{lemma}[Restriction to a convex cap]\label{j:convex-cap}
Let $K\subset D$ be finite convex tessellation polygons sharing a
starting port with the same inward direction. Let $L$ be the raw weight
of first-exit chords of $D$ other than those confined to $K$ and ending
at a boundary port shared by $K$ and $D$. Let $E$ be
the raw weight of first-exit chords of $K$ ending on ports newly created
by the restriction. Then
\[
 \sin\alpha\,L\le E\le L/\sin\alpha.
\]
The same conclusion holds with $D$ the saturated tessellation
half-plane and $K$ a finite convex cap.
\end{lemma}

\begin{proof}
The common chords, confined to $K$ and ending on a shared true boundary
port, have the same intrinsic turning phase in both balances and cancel.
The two remaining complex sums are equal. Every real coefficient in
either convex domain lies between $\sin\alpha$ and one. Taking real
parts proves both inequalities. For the half-plane, use the convex
exhaustion justified by saturation in the wedge argument.
\end{proof}

Recall the regular tessellation hexagon $D_R$, the ordered first-exit
chords $\mathcal C_R$ of diameter at least $R/10$, and their raw
partition $Z_R^{\rm chord}$. For an upper bound at a source $a$, choose
a small similarly aligned tessellation hexagon $P$ contained within
distance $R/25$ of $a$, with
$\operatorname{dist}(a,\partial P)\ge cR$. Such a choice is possible
with fixed constants for large $R$ by taking its center within a bounded
lattice distance of $a$. A chord of diameter at least $R/10$ reaches
distance at least $R/20$ from $a$, and hence is lost on restriction to
$D_R\cap P$. Lemma~\ref{j:convex-cap} bounds its weight by a constant
times the artificial exits of that intersection. These exits are a
subset of the artificial exits of $P$ intersected with the supporting
tessellation half-plane at $a$. Applying the lemma to this half-plane
cap bounds them by $CW_{cR}$. Thus the macroscopic weight at every
source is at most $R^{-1/4+o(1)}$.

For a lower bound choose the ports in a fixed middle portion of one
side, at distance at least $cR>R/10$ from every other side. A half-plane
arch from such a source reaching radius $R$ is either a macroscopic
chord contained in $D_R$ or is lost on restriction to $D_R$. By
Lemma~\ref{j:convex-cap}, all lost arch weight is bounded by a constant
times the exits of $D_R$ on the other sides, and these are macroscopic
as well. Thus the macroscopic weight at each of these sources is at
least a constant times $W_R=R^{-1/4+o(1)}$. There are $O(R)$ boundary
ports in all and order $R$ chosen ports. We have proved
\begin{equation}\label{j:chord-partition}
 Z_R^{\rm chord}=R^{3/4+o(1)}.
\end{equation}

\begin{proposition}[Interior visits and nested polygons]\label{j:visit}
Let $D$ be a finite convex tessellation polygon and let $v$ be a
honeycomb vertex not adjacent to a boundary port. Sum over all ordered
first-exit boundary chords of $D$ that visit $v$, and put
\[
 S_D(v)=\sum_{\gamma\ni v}x^{N(\gamma)}\cos(3T(\gamma)/8).
\]
Let $Z_v$ be the partition of disjoint unoriented cycles in $D$
strictly enclosing $v$ and avoiding it, with cycle weight
$2x^{|\lambda|}$. Let $Z_{\rm out}(\lambda)$ count disjoint cycles in
$D$ surrounding and disjoint from $\lambda$ with the same weights. At a
vertex of $\lambda$, call the corner convex or reflex when its interior
angle is $2\pi/3$ or $4\pi/3$, respectively, and put
\[
 c_\lambda=
 \begin{cases}
 \sqrt2,&v\text{ is a convex corner of }\lambda,\\
 2+\sqrt2,&v\text{ is a reflex corner of }\lambda.
 \end{cases}
\]
Then
\begin{equation}\label{j:visit-identity}
 S_D(v)=3Z_v+\sum_{\lambda\ni v}c_\lambda
       x^{|\lambda|}Z_{\rm out}(\lambda).
\end{equation}

For each of the three incident face centers $f$, let $Z_D(f)$ be the
partition of disjoint cycles in $D$ surrounding $f$, again with weight
$2x^{|\lambda|}$. Put $k_\lambda=1$ at a convex corner and
$k_\lambda=2$ at a reflex corner. Then
\begin{equation}\label{j:incident-face-partition}
 \sum_{f\sim v}Z_D(f)
 =3Z_v+\sum_{\lambda\ni v}2k_\lambda x^{|\lambda|}
                         Z_{\rm out}(\lambda).
\end{equation}
In particular the ordinary visit mass
$\sum_{\gamma\ni v}x^{N(\gamma)}$ is comparable, with absolute positive
constants, to $\sum_{f\sim v}Z_D(f)$.
\end{proposition}

\begin{proof}
Remove the triangle centered at $v$, creating three inner ports. For
each outer starting port, the signed balance still holds in the
punctured domain: its source is outside every possible first-collision
loop, including a loop surrounding the removed triangle, so the same
collision cancellation applies. Subtract the punctured balance from
the original one. The outer chords avoiding $v$ cancel, leaving the
outer chords through $v$ equal to the paths ending at an inner port.
Take real parts, sum over the outer sources, and reverse these paths.
Reversal changes $T$ to $-T$ and preserves the real part. Thus $S_D(v)$
is the real sum of paths from the three inner ports to the outer ports.

For each inner source, attach an arbitrary configuration counted by
$Z_v$ as a background and grow the weighted path tree in the punctured
domain. Stop at an exit, a first self-collision, or first contact with a
background cycle. The initial total strength of the three trees is
$3Z_v$. A self-collision loop not enclosing $v$ is approached from
outside and cancels with its reverse exactly as in
Section~\ref{j:local}. A loop enclosing $v$ is approached from inside,
because the inner source is in the same component as $v$. Its corner at
the collision vertex is reflex: the other loop vertices contribute
$\pm7\pi/3$ of turning, and departure from the incoming stem contributes
$\pm\pi/3$. The two orientations of that loop therefore contribute
\[
 2x^{|\lambda|}
 \cos\left(\frac38\frac{8\pi}{3}\right)=-2x^{|\lambda|}.
\]
This cancels first contact with the same loop added to the background.
Conversely, remove the first contacted, necessarily innermost,
background cycle and follow it from the contact in either orientation.
First contact ensures that this cycle has no earlier intersection with
the stem. This is the inverse operation, so every background-contact
term is canceled exactly once.

An outer exit can coexist with no background cycle, since the source
starts inside every such cycle. Its real sum is the $S_D(v)$ identified
above. An exit to another inner port, closed through $v$, gives a simple
cycle $\lambda$ through $v$, and every remaining background cycle
surrounds it. If $L=|\lambda|$, opening the cycle at $v$ leaves two
directed inner-return paths of weight $x^{L-1}$ each. At a convex corner
their turns are $\pm5\pi/3$, and at a reflex corner they are
$\pm7\pi/3$: the full loop turn is $\pm2\pi$, with the turn at $v$
removed. Write $t_\lambda=5\pi/3$ or $7\pi/3$ in these two cases.
The remaining real tree identity is
\[
 3Z_v=S_D(v)+
 \sum_{\lambda\ni v}2x^{|\lambda|-1}
 \cos\left(\frac38\,t_\lambda\right)Z_{\rm out}(\lambda).
\]
Moving the inner-return terms to the other side supplies the required
factor $x^{-1}$:
\[
 -\frac2x\cos(5\pi/8)=\sqrt2,\qquad
 -\frac2x\cos(7\pi/8)=2+\sqrt2.
\]
This proves \eqref{j:visit-identity}.

A cycle avoiding $v$ places all three incident face centers in the same
component as $v$, because it uses none of the three edges incident to
$v$. A cycle through $v$ encloses exactly one incident face center at
a convex corner and two at a reflex corner. In a surrounding
configuration it must be the innermost cycle: an inner cycle avoiding
$v$ would also enclose $v$, which lies on the outer cycle. These facts
partition the configurations counted by the three $Z_D(f)$ and prove
\eqref{j:incident-face-partition}. Comparing the positive coefficients
in the two identities gives
\[
 \frac1{\sqrt2}\sum_{f\sim v}Z_D(f)\le S_D(v)
 \le\sum_{f\sim v}Z_D(f).
\]
Finally convex boundary positivity gives
\[
 S_D(v)\le\sum_{\gamma\ni v}x^{N(\gamma)}
 \le\frac{S_D(v)}{\sin\alpha},
\]
which proves the ordinary visit comparison.
\end{proof}

Let $\mathcal Z(R)$ be the $+2$ partition around the central face center
in $D_R$. For a honeycomb vertex $v$ in $D_R$, write
\[
 \mathcal V_R(v)=\sum_{\substack{\gamma\in\mathcal C_R\\v\in\gamma}}
                  x^{N(\gamma)}.
\]
This is a raw mass under the finite chord ensemble, so
\[
 \mathbb E_R^{\rm chord}N
 =\frac1{Z_R^{\rm chord}}\sum_{v\in D_R}\mathcal V_R(v).
\]
For $v$ at distance at most $R/10$ from the center, every first-exit
boundary chord visiting $v$ is macroscopic for large $R$. A translated
hexagon of side $\lfloor R/2\rfloor$ centered at each incident face lies
inside $D_R$. Positivity under domain inclusion and
Proposition~\ref{j:visit} therefore give
\[
 \mathcal V_R(v)\ge c\mathcal Z(\lfloor R/2\rfloor)
 \quad\text{uniformly for these bulk vertices}.
\]
For every vertex not adjacent to a boundary port, a translated hexagon
of side $3R$ centered at an incident face contains $D_R$. The same
proposition gives
$\mathcal V_R(v)\le C\mathcal Z(3R)$ uniformly. There are order $R^2$
bulk vertices and at most order $R$ vertices adjacent to boundary
ports. Each of the latter contributes at most one to the normalized
mean, since a chord visits a vertex at most once. Hence
\[
 \begin{aligned}
 \frac{cR^2\mathcal Z(\lfloor R/2\rfloor)}{Z_R^{\rm chord}}
 &\le\mathbb E_R^{\rm chord}N\\
 &\le\frac{CR^2\mathcal Z(3R)}{Z_R^{\rm chord}}+CR.
 \end{aligned}
\]
Using \eqref{j:chord-partition} and $\mathcal Z\ge1$ yields
\[
 \begin{aligned}
 R^{5/4+o(1)}\mathcal Z(\lfloor R/2\rfloor)
 &\lesssim \mathbb E_R^{\rm chord}N\\
 &\lesssim R^{5/4+o(1)}\mathcal Z(3R).
 \end{aligned}
\]
In particular $\mathbb E_R^{\rm chord}N\ge R^{5/4+o(1)}$.
Section~\ref{c7:localization} will supply
$\mathcal Z(R)=R^{1/12+o(1)}$, completing the mean exponent $4/3$ for
this finite chord law.
